\documentclass[book.tex]{subfiles}

\begin{document}

\chapter{Schrodinger's Equation}

\label{chap:schrodinger}
\noindent\rule{11cm}{0.4pt} \\
\textbf{Prerequisites:} \autoref{chap:quantum_mechanics}, \autoref{chap:position_and_momentum} \\
\textbf{Difficulty Level:} ** \\
\noindent\rule{11cm}{0.4pt}
\vspace{5mm}

We previously learned about Newton's laws, and Maxwell's equations which respectively describe how the time evolution of classical mechanical and electromagnetic systems proceeded. Schrodinger's equation provides an analogous framework for describing the evolution of quantum mechanical systems.

\section{Basic Formulation}

Let $\psi(t)$ be the wave function of a quantum mechanical system at time $t$. Then the wave function evolves over time according to Schrodinger's equation, written as 
\begin{equation}
i \hbar \frac{d}{dt}|\psi(t) \rangle = \hat{H}|\psi(t) \rangle
\end{equation}
where $\hat{H}$ is an operator, called the \textit{Hamiltonian}, on the space of wave functions. Given the ubiquity of the Hamiltonian in quantum mechanics we will also often leave off the hat and use a plain $H$ for the Hamiltonian.

What is the Hamiltonian? We described that the process of quantization converts classical vectorial or scalar quantitites into operators. One way of thinking about the Hamiltonian is as the quantization of the classical energy. Recall quantization involves swapping out scalar quantities for operators, so we describe the type of the Hamiltonian as

\begin{align}
H: \mathcal{H} \to \mathcal{H}
\end{align}

where $\mathcal{H}$ is the Hilbert space of wave functions. We use $\mathcal{H}$ in place of concrete Hilbert spaces like $L^2(\mathbb{R})$ or $L^2(\mathbb{R}^3)$ since the Hamiltonian formalism holds true for all of these specific choices. The intuition of Hamiltonian as the quantized energy drives the definition of many simple Hamiltonians for quantum systems. Let's look at an example of the usual form of a Hamiltonian.

\begin{equation}
    \hat{H} = \frac{\hat{p}^2}{2m} + \hat{V} 
\end{equation}
Here $\frac{\hat{p}^2}{2m}$ can be thought of as the kinetic energy operator and $\hat{V}$ as the potential energy operator.

This equation should look familiar to you. If we undo the quantization, this converts into the classical energy equation
\begin{equation}
E = \frac{p^2}{2m} + V(x)
\end{equation}

We can express the Hamiltonian in Physika for the concrete Hilbert space $L^2(\mathbb{R})$ as 
\begin{lstlisting}[language=python]
m : $\mathbb{R}^+$, V: $L^2(\mathbb{R})$
def $\hat{H}$($\psi$: $L^2(\mathbb{R})$) $\to$ $L^2(\mathbb{R})$:
  return $\hat{p}^2$($\psi$)/2m + V * $\psi$
\end{lstlisting}

It is important to note that the choice of Hamiltonian $\bar{H}$ is an act of \textit{modeling}. That is, while the Hamiltonian should generally correspond to the energy for a system, in practice complex systems often require simplifications or tweaks. When faced with a new physical system, there is not one canonical procedure for choosing a Hamiltonian for a physical system, but the physicist can derive a Hamiltonian that appears to model the system from first principles or from known Hamiltonian for simpler systems. As we will see later in this book, it is even possible to learn an approximate Hamiltonian given sufficient data about the evolution of a quantum system.

Let us now return to Schrodinger's equation. Rewrite it as follows
\begin{equation}
    \frac{\partial}{\partial t} | \psi \rangle = -i \hbar \hat{H} | \psi(t) \rangle 
\end{equation}
This equation is roughly similar to an equation for gradient descent. Very approximately the Schrodinger's equation describes a form of gradient descent, by compelling the wave function to evolve in a way favorable to the energy of the system.

\section{Forces from Schrodinger's Equation}

One qualitative difference in quantum versus classical mechanics is that we have less reason to work directly with forces. We note that it is possible to define a force operator
\begin{align}
\hat{F} &= \frac{d}{dt} \hat{p}
\end{align}
However, it turns out that the more useful quantity is given by
\begin{align}
    \frac{d}{dt} \langle \hat {p} \rangle
\end{align}
Ehrenfest's theorem shows that the expectation value of the position operator evolves in a direction controlled by this modified quantity
\begin{align}
    m \frac{d}{dt} \langle \hat{x} \rangle &= \langle \hat{p} \rangle  \\
    \frac{d}{dt} \langle \hat{p} \rangle &= - \langle V'(\hat{x}) \rangle
\end{align}

\end{document}
